\begin{theorem}[Completely regular space lattice dependency route]
\label{thm:completely-regular-space-lattice-dependency-route}
Every $\CompletelyRegularSpaceUp$ separator read first consumes the $\TopologyUp$ carrier of \autoref{ch:concrete-instances-topology-namecert} and the $\RegularTopologicalSpaceUp$ point-versus-closed-set source of \autoref{ch:concrete-instances-regular-topological-space-namecert}. Only after those rows are displayed may the packet export its $\RealUp$-valued $\ContinuousMapUp$ separator row to $\UrysohnLemmaUp$ or $\UrysohnMetrizationUp$ consumers. The route is the finite lattice edge $T,R,C,F,H,K,P,N$ of \autoref{def:completely-regular-space-carrier}, not a host powerset, selected complement, arbitrary function space, normality theorem, or metrization proof.
\end{theorem}
\begin{proof}
Project the carrier $Q=(T,R,C,F,H,K,P,N)$. The topology row $T$ supplies the displayed open and closed vocabulary, while the regular topological row $R$ binds the point-versus-closed-set request to that same source.

The request row $C$ must be visible before the separator row $F$ is exported. Since $F$ is a $\RealUp$-valued $\ContinuousMapUp$ row, an Urysohn-facing consumer receives a displayed separator rather than an ambient family of functions.

Finally the structural rows $H,K,P,N$ preserve the same route for $\UrysohnLemmaUp$ and $\UrysohnMetrizationUp$ consumers. No carrier coordinate supplies a hidden powerset, selected complement, normality theorem, or metrization proof, so the lattice edge is exactly the displayed topological dependency route.
\end{proof}
